% preamble-firstwords.tex
% -----------------------------------------------------------------------
% Lo propio de «First Words» (firstwords.tex, firstwords-bw.tex), el
% cuaderno de primeras palabras en inglés, cargado DESPUÉS de preamble.tex
% y de preamble-palabras.tex: la página, las cajas de actividad y el pie
% son los del cuaderno de primeras palabras en español (ver
% notes/06-first-words.md); aquí solo cambia la tarjeta de la palabra.
%
% En español, la tarjeta enseña la palabra partida en sílabas y entera.
% En inglés se lee por sonidos, así que la tarjeta enseña la palabra
% entera con un "botón" debajo de cada sonido (sound buttons, como en
% los colegios ingleses): un punto si el sonido es una sola letra (c, a,
% t), una raya si son varias (sh, ee, ck), y un arco que une la vocal con
% la e del final si es una e mágica (cake). Se lee tocando cada botón y
% diciendo su sonido, y luego la palabra de corrido.
% -----------------------------------------------------------------------

\usetikzlibrary{calc}

% Sin ligaduras: Andika junta ff, fi y fl en una sola letra, con una sola
% barra, y en este cuaderno cada letra cuenta -- off se lee o-ff, y esas
% dos efes tienen que verse como dos, igual que en la tabla de sonidos de
% cualquier colegio. La misma letra de preamble.tex, solo sin "liga" (los
% guiones y las comillas de TeX, -- y ``, siguen funcionando).
\setmainfont{Andika}[
  Path = fonts/andika/,
  Extension = .ttf,
  UprightFont = *-Regular,
  BoldFont = *-Bold,
  ItalicFont = *-Italic,
  BoldItalicFont = *-BoldItalic,
  Ligatures = NoCommon,
]
\setsansfont{Andika}[
  Path = fonts/andika/,
  Extension = .ttf,
  UprightFont = *-Regular,
  BoldFont = *-Bold,
  ItalicFont = *-Italic,
  BoldItalicFont = *-BoldItalic,
  Ligatures = NoCommon,
]

% Cada grafema es un nodo pegado al anterior por la línea base
% (tools/gen_palabras.py, tarjeta_sonidos): la palabra se ve entera, como
% si fuera un solo texto, y los botones caen debajo de sus letras.
\tikzset{grafema/.style={anchor=base west, inner sep=0pt, outer sep=0pt}}

% Los botones, en unidades de la letra (em), para que sirvan igual con la
% palabra enorme del otoño que con las tres más pequeñas de la
% primavera. A 0,36 em por debajo de la línea base quedan por debajo de
% las letras que bajan (p, g, y).
\newcommand{\botonPunto}[1]{%
  \fill[colorLectura] ($(#1.base)+(0,-0.36em)$) circle (0.075em);}
\newcommand{\botonRaya}[1]{%
  \draw[colorLectura, line width=0.06em, line cap=round]
    ($(#1.base west)+(0.12em,-0.36em)$) -- ($(#1.base east)+(-0.12em,-0.36em)$);}
% La e mágica: un arco por debajo, de la vocal a la e, que pasa por
% debajo del botón de la consonante de en medio.
\newcommand{\botonArco}[2]{%
  \draw[colorLectura, line width=0.06em, line cap=round]
    ($(#1.base)+(0,-0.3em)$) .. controls +(0.05em,-0.42em) and +(-0.05em,-0.42em)
    .. ($(#2.base)+(0,-0.3em)$);}

% La tarjeta: #1 el trimestre (el tamaño, \fuentePalabra), #2 los nodos
% y los botones. La letra se elige antes del dibujo para que los em de
% los botones sean los de la palabra.
\newcommand{\tarjetaSonidos}[2]{%
  {\fuentePalabra{#1}\ajustado{\begin{tikzpicture}[baseline=(g1.base)]#2\end{tikzpicture}}\par}%
}

% Un nombre que se lee de un golpe (Toby, Amy): la palabra entera, sin
% botones, con el mismo aire debajo que una tarjeta con botones.
\newcommand{\tarjetaEntera}[2]{%
  {\fuentePalabra{#1}\ajustado{#2}\par\vspace{0.3em}}%
}

% Una tricky word, los viernes de primavera (tools/gen_palabras.py,
% lectura_tricky): con su casilla, como las palabras de la semana, pero
% dentro de un marco y sin botones -- se lee entera.
\newcommand{\palabraTricky}[1]{\mbox{\casillaGrande\hspace{2mm}%
  \fcolorbox{colorCrea}{white}{\fontsize{24}{28}\selectfont\strut\ #1\ }}}

% =========================================================================
% La tabla de sonidos (frontmatter/firstwords/sonidos.tex; las secciones
% las genera tools/gen_palabras.py en content/firstwords/generated-sonidos.tex)
% =========================================================================
% Cada sonido en grande, y debajo una palabra que lo tiene, con sus
% letras resaltadas (en blanco y negro, la negrita sola ya lo dice). Las
% casillas son cajas de ancho fijo dentro de un párrafo: salen nueve por
% línea -- los nueve sonidos de la segunda parte del otoño, una línea
% justa --, y el párrafo las parte donde haga falta. El \strut iguala el
% alto y el fondo de todos los sonidos, bajen o no de la línea (g, j, p,
% y): así las palabras de debajo quedan todas a la misma altura.
\newcommand{\resalta}[1]{\textcolor{colorLectura}{\bfseries #1}}
\newcommand{\seccionSonidos}[1]{%
  \par\vspace{4mm}{\large\bfseries #1\par}\vspace{1mm}%
  \noindent\raggedright\ignorespaces}
\newcommand{\sonido}[2]{%
  \makebox[18mm][c]{\begin{tabular}[t]{@{}c@{}}%
    {\fontsize{26}{30}\selectfont\strut #1}\\[-0.8mm]{\large #2}%
  \end{tabular}}\hspace{0.8mm plus 1mm}\ignorespaces}
\newcommand{\finSonidos}{\par}

